\fancyfoot[R]{v1.0.1}
\section{Pseudoinverse matrices. Skeletonization. Singular value decomposition (SVD)}

\begin{problem}{Skeletonization and pseudoinverse matrix}
    Find the pseudoinverse matrix to matrix $A$ using skeletonization
    \[
             A = \begin{bmatrix}
                    2 & -1 & 0\\
                    -1 & 1 & 1\\
                    0 & 1 & 2
                \end{bmatrix}.
    \]
\end{problem}
\begin{solution}
    Let's start with skeletonization and then find the pseudoinverse matrix by the formula
    \[
        A^+ = G^*(GG^*)^{-1}(F^*F)^{-1}F^*, 
    \]
    \[
        A = \Vast[
            \overbrace{\fbox{$\begin{matrix}
                2 & -1 \\
                -1 & 1 \\
                0 & 1
            \end{matrix}$}}^{F} \hspace*{0.25cm} \begin{matrix}
                0 \\
                1 \\
                2
            \end{matrix}\Vast] \Rightarrow \begin{bmatrix}
                2 & -1 & 0 \\
                0 & \medmath{\frac{1}{2}} & 1 \\
                0 & 1 & 2
            \end{bmatrix} \Rightarrow \begin{bmatrix}
                1 & -\medmath{\frac{1}{2}} & 0\\
                0 & \medmath{\frac{1}{2}} & 1\\
                0 & 0 & 0
            \end{bmatrix} \Rightarrow \Vast[
                \begin{array}{c}
                    \overbrace{\fbox{$\begin{matrix}
                        1 & 0 & 1\\
                        0 & 1 & 2
                    \end{matrix}$}}^{G}\\
                    \hspace*{0.18cm} \begin{matrix}
                        0 & 0 & 0\\\\
                    \end{matrix}
                \end{array}
                \Vast].
    \]
    Let's check
    \[
        F \cdot G = \begin{bmatrix}
            2 & -1 \\
            -1 & 1 \\
            0 & 1
        \end{bmatrix} \cdot \begin{bmatrix}
            1 & 0 & 1 \\
            0 & 1 & 2
        \end{bmatrix} = \begin{bmatrix}
            2 & -1 & 0\\
            -1 & 1 & 1\\
            0 & 1 & 2
        \end{bmatrix}.
    \]
    Correct. Now we need to find the pseudoinverse matrix
    \[
        G^*(G G^*)^{-1} = \begin{bmatrix}
            1 & 0 \\
            0 & 1 \\
            1 & 2
        \end{bmatrix} \cdot \left(\begin{bmatrix}
            1 & 0 & 1\\
            0 & 1 & 2
        \end{bmatrix} \cdot \begin{bmatrix}
            1 & 0 \\
            0 & 1\\
            1 & 2
        \end{bmatrix}\right)^{-1} = \begin{bmatrix}
            1 & 0\\
            0 & 1\\
            1 & 2
        \end{bmatrix} \cdot \dfrac{1}{6} \begin{bmatrix}
            5 & -2 \\
            -2 & 2
        \end{bmatrix} = \dfrac{1}{6}\cdot\begin{bmatrix}
            5 & -2 \\
            -2 & 2\\
            1 & 2
        \end{bmatrix},  
    \]
    \begin{eqnarray*}
        (F^* F)^{-1}  F^* &=& \left(\begin{bmatrix}
            2 & -1 & 0\\
            -1 & 1 & 1
        \end{bmatrix} \cdot \begin{bmatrix}
            2 & -1 \\ -1 & 1 \\
            0 & 1
        \end{bmatrix}\right)^{-1}\cdot \begin{bmatrix}
            2 & -1 & 0\\
            -1 & 1 & 0
        \end{bmatrix} = \begin{bmatrix}
            5 & -3 \\
            -3 & 3
        \end{bmatrix}^{-1} \begin{bmatrix}
            2 & -1 & 0\\
            -1 & 1 & 0
        \end{bmatrix} = \\ &=& \dfrac{1}{6} \begin{bmatrix}
            3 & 3 \\
            3 & 5
        \end{bmatrix} \cdot \begin{bmatrix}
            2 & -1 & 0\\
            -1 & 1 & 0
        \end{bmatrix} = \dfrac{1}{6} \begin{bmatrix}
            3 & 0 & 0\\
            1 & 2 & 0
        \end{bmatrix}.
    \end{eqnarray*}
    All things considered, we can obtain pseudoinverse matrix $A^+$
    \[
        A^+=\dfrac{1}{36} \begin{bmatrix}
            5 & -2 \\
            -2 & 2\\
            1 & 2
        \end{bmatrix} \cdot \begin{bmatrix}
            3 & 0 & 0\\
            1 & 2 & 0
        \end{bmatrix} = \dfrac{1}{36} \begin{bmatrix}
            13 & -4 & 0\\
            -4 & 4 & 0\\
            5 & 4 & 0
        \end{bmatrix}.
    \]
\end{solution}
    \begin{note}{}{}
        Matrix $(G G^*)$ is called Gramm matrix and contains results of scalar products. 
    \end{note}

    \begin{definition}{SVD}{}
        
    Singular Value Decomposition of matrix $A\in M_{m\times n}(\C)$ is a decomposition of a kind
    \[
        A = U\Sigma V^*,  
    \]
    where $U \in M_{m\times m}(\C)$ and $V \in M_{n\times n}(\C)$ are unitary matrices
    \begin{gather*}
        U^*U = UU^*=UU^{-1} = I\\
        V^*V = VV^*=VV^{-1} = I
    \end{gather*}
    and $\Sigma \in M_{m\times n}(\C)$ is a diagonal matrix of a kind
    \[
        \scalebox{1.5}{$\Sigma = $}
        \left[\begin{tabular}{clp{0.1cm}p{0.1cm}}
            \multicolumn{1}{l}{\scalebox{1.25}{$\sigma_1$}}        &         & & \multicolumn{1}{l}{\multirow{3}{*}{\begin{tikzpicture}[x=0.55pt,y=0.55pt,yscale=-1,xscale=1]
                    \draw  [line width=1.75]  (304.32,138.5) .. controls (312.66,138.5) and (319.43,147.68) .. (319.43,159) .. controls (319.43,170.32) and (312.66,179.5) .. (304.32,179.5) .. controls (295.98,179.5) and (289.21,170.32) .. (289.21,159) .. controls (289.21,147.68) and (295.98,138.5) .. (304.32,138.5) -- cycle ;
                    \end{tikzpicture}}} \\ %! 
\multicolumn{1}{l}{}         & \scalebox{1.25}{$\sigma_2$}       & \multicolumn{1}{c}{} & \multicolumn{1}{c}{}                     \\ %! 
            \multicolumn{2}{c}{\multirow{2}{*}{
                \begin{tikzpicture}[x=0.55pt,y=0.55pt,yscale=-1,xscale=1]
                    \draw  [line width=1.75]  (304.32,138.5) .. controls (312.66,138.5) and (319.43,147.68) .. (319.43,159) .. controls (319.43,170.32) and (312.66,179.5) .. (304.32,179.5) .. controls (295.98,179.5) and (289.21,170.32) .. (289.21,159) .. controls (289.21,147.68) and (295.98,138.5) .. (304.32,138.5) -- cycle ;
                    \end{tikzpicture}
            }} & \scalebox{1.25}{$\ddots$}              & \multicolumn{1}{l}{}  \\ %!
            \multicolumn{2}{c}{}                   &                       & \multicolumn{1}{l}{\scalebox{1.25}{$\sigma_r$}}   \\[0.05cm] %!
 
            \multicolumn{4}{c}{\begin{tikzpicture}[x=0.55pt,y=0.55pt,yscale=-1,xscale=1]
                \draw [line width=1.75]  (246,143.5) .. controls (246,133.56) and (273.41,125.5) .. (307.21,125.5) .. controls (341.02,125.5) and (368.43,133.56) .. (368.43,143.5) .. controls (368.43,153.44) and (341.02,161.5) .. (307.21,161.5) .. controls (273.41,161.5) and (246,153.44) .. (246,143.5) -- cycle ;
                \end{tikzpicture}
                }                                                                        
            \end{tabular}\right],
    \]
    where $\sigma_1 \geq \sigma_2 \geq \ldots \geq \sigma_r$ -- singular values and $r = \rank A$.
    \end{definition}
    \begin{theorema}{How to find SVD}{}
        The eigenvectors of $A^*A$ make up the columns of $V$ , the eigenvectors of $AA^*$  make up the columns of $U$, the singular values in $\Sigma$ are square roots of eigenvalues from  $AA^*$ or  $A^*A$. 
    \end{theorema}
    
    \begin{theorema}{How to find pseudoinverse using SVD}{}
        Let $U\Sigma V^*$ be the SVD of a matrix $A\in M_{m\times n}(\C)$, then
        $$
            A^+=V\Sigma^+ U^*.
        $$
    \end{theorema}

T\begin{problem}{Singular Value Decomposition}
Find singular value decomposition of a matrix
\[
    A = \begin{bmatrix}
        3 & 3 \\
        0 & 0 \\
        4 & 4
    \end{bmatrix}.  
\]
\end{problem}
\begin{solution}
    Let's start with finding $\Sigma \in M_{2\times 2}(\R)$
    \[
        A^* A = \begin{bmatrix}
            3 & 0 & 4\\
            3 & 0 & 4
        \end{bmatrix} \cdot \begin{bmatrix}
            3 & 3\\
            0 & 0\\
            4 & 4
        \end{bmatrix} = \begin{bmatrix}
            25 & 25 \\ 
            25 & 25
        \end{bmatrix}.   
    \]
    Next step is to find eigenvalues and eigenvectors for the matrix $A^* A$.
    \[
        (A^* A - \lambda_2 I) \vec v_2=0
        \Leftrightarrow
        \det \left(\begin{bmatrix}
            25-\lambda & 25 \\
            25 & 25-\lambda
        \end{bmatrix}\right) = 0 
        \Leftrightarrow \lambda^2 - 50\lambda = 0 
        \Leftrightarrow \lambda_1 = 50, \ \lambda_2 = 0.
    \]
    Hence singular values are equal to $\sigma_1 = \sqrt{50}, \ \sigma_2 = 0$ and 
    \[
        \scalebox{1.5}{$\Sigma$} = \begin{bmatrix}
            \sqrt{50} & 0\\
            0 & 0 \\
            0 & 0
        \end{bmatrix}.  
    \]
    Let's  find  $V \in M_{2\times 2}(\R)$. We calculate eigenvalues of previously  obtained matrix $A^* A$.
    \[
        (A^* A - \lambda_1 I) \vec v_1 = 0
        \Leftrightarrow
        \begin{bmatrix}
            -25 & 25 \\
            25 & -25
        \end{bmatrix} \vec v_1 = 0.
    \]
    One of the solutions is  $\vec v_1 = \begin{bmatrix}
        1 \\ 1
    \end{bmatrix}$. We need to normalize it, so $\vec v_1 = \dfrac{\sqrt{2}}{2}\begin{bmatrix}
        1 \\ 1
    \end{bmatrix}$. Similarly by solving $(A^* A - \lambda_2 I) \vec v_2-0$, we get $\vec v_2 = \dfrac{\sqrt 2}{2}\begin{bmatrix}
        -1 \\ 1
    \end{bmatrix} $.  Combining  $\vec v_1$ and $\vec v_2$, we get
    \[
        V = \dfrac{\sqrt{2}}{2} \begin{bmatrix}
            1 & -1\\ 1 & 1
        \end{bmatrix}.  
    \] 
    Let's find matrix $U \in M_{3\times 3}(\R)$, which is constructed of a new matrix $AA^*$ eigenvectors. 
    \[
        AA^* = \begin{bmatrix}
            3 & 3 \\
            0 & 0\\
            4 & 4
        \end{bmatrix} \cdot \begin{bmatrix}
            3 & 0 & 4\\
            3 & 0 & 4
        \end{bmatrix} = \begin{bmatrix}
            18 & 0 & 24\\
            0 & 0 & 0\\
            24 & 0 & 32
        \end{bmatrix}.
    \]
    We calculate the new matrix $AA^*$  eigenvectors
    \[      
        (AA^*-\lambda_1 I)\vec u_1 = 0
        \Leftrightarrow
        \begin{bmatrix}
            -32 & 0 & 24\\
            0 & -50 & 0\\
            24 & 0 & -18
        \end{bmatrix}\vec u_1 = 0.
        \]
    Solving by Gaussian elimination
    \[
        \begin{bmatrix}[ccc|c]
            -4 & 0 & 3 & 0\\
            0 & -1 & 0 & 0\\
            4 & 0 & -3 & 0
        \end{bmatrix}  \sim \begin{bmatrix}[ccc|c]
            -1 & 0 & \medmath{\frac{3}{4}} & 0\\
            0 & 1 & 0 & 0\\
            0 & 0 & 0 & 0
        \end{bmatrix} \Rightarrow \vec u_1 = \dfrac{1}{5}\begin{bmatrix}
        \medmath{3} \\[0.2cm] 0 \\ 4
    \end{bmatrix}.
    \]
    Similarly by solving $(AA^*-\lambda_2 I)\vec u=0$,  we get
    \[
        \begin{bmatrix}[ccc|c]
            18 & 0 & 24 & 0\\
            0 & 0 & 0 & 0\\
            24 & 0 & 32 & 0
        \end{bmatrix}  \sim \begin{bmatrix}[ccc|c]
            1 & 0 & \medmath{\frac{4}{3}} & 0\\
            0 & 0 & 0 & 0\\
            0 & 0 & 0 & 0
        \end{bmatrix} \Rightarrow 
        \vec u_2 = 
        \begin{bmatrix}
            0 \\ 1 \\ 0
        \end{bmatrix},
        \vec u_3 = 
        \dfrac{1}{5}
        \begin{bmatrix}
            -4 \\ 0 \\ 3
        \end{bmatrix}.
    \]
    All things considered, the result of singular value decomposition is the following
    $$
    A=U\Sigma V^*,
  $$
  where 
    \[
        U = \dfrac{1}{5}\begin{bmatrix}
            3 & 0 & -4\\
            0 & 5 & 0\\
            4 & 0 & 3
        \end{bmatrix}, \vspace*{0.5cm} \Sigma = \begin{bmatrix}
            \sqrt{50} & 0\\
            0 & 0\\
            0 & 0
        \end{bmatrix}, \vspace*{0.5cm} V = \frac{\sqrt{2}}{2} \begin{bmatrix}
            1 & -1 \\
            1 & 1
        \end{bmatrix}.
    \]
\end{solution}
\begin{note}
    We could find more simple decomposition of the matrix $A$ by omitting all calculations with zero singular values. Such decomposition is called compact version of the SVD
    $$
        A=U_{r}\Sigma_{r} V_{r}^*,\quad r=\rank A.
    $$
    For example, matrix $A$ from the previous problem has following compact SVD
    $$
        A=U_{r}\Sigma_{r} V_{r}^*,
    $$
    where
    \[
        U_{r} = \dfrac{1}{5}\begin{bmatrix}
            3 \\
            0 \\
            4  
        \end{bmatrix}
            , \vspace*{0.5cm} \Sigma_{r} = 
        \begin{bmatrix}
            \sqrt{50}
        \end{bmatrix}
            , \vspace*{0.5cm} V_{r} = \frac{\sqrt{2}}{2} 
        \begin{bmatrix}
            1 \\
            1 
        \end{bmatrix}.
    \]
\end{note}
